% „Das kann ich“ (zusammenbau v0.9): je Kette eine Zeile – Zone nur im Gesamt (\mitzone)
\begin{lbdaskannich}
\ifdefined\mitzone
\abhakgruppe{Das kennst du schon}
\abhak{1, 5}{Ich kann negative Zahlen addieren und subtrahieren.}
\abhak{2}{Ich kann Zahlen mit Vorzeichen malnehmen.}
\abhak{3}{Ich kann mit Dezimalzahlen und einfachen Brüchen rechnen.}
\abhak{4}{Ich kann Punkt vor Strich rechnen.}
\fi
\abhakgruppe{Einheit 1 · Terme aufstellen und berechnen}
\abhak{6–8, 11–14}{Ich kann einen Term aufstellen.}
\abhak{9}{Ich kann Termwerte berechnen.}
\abhak{10}{Ich kann zu einem Term eine passende Situation erfinden.}
\abhakgruppe{Einheit 2 · Terme zusammenfassen}
\abhak{15}{Ich erkenne gleichartige Glieder.}
\abhak{16}{Ich kann die Vorzahl eines Gliedes ablesen.}
\abhak{17}{Ich kann die Glieder eines Terms mit ihrem Vorzeichen aufschreiben.}
\abhak{18–20, 23–26}{Ich kann Terme zusammenfassen.}
\abhak{21–22}{Ich kann Terme malnehmen.}
\abhakgruppe{Einheit 3 · Klammern auflösen}
\abhak{27}{Ich erkenne, was vor einer Klammer steht.}
\abhak{28–34}{Ich kann Klammern auflösen.}
\abhakgruppe{Einheit 4 · Ausklammern}
\abhak{35}{Ich erkenne, welcher Faktor in jedem Glied steckt.}
\abhak{36–43}{Ich kann ausklammern.}
\end{lbdaskannich}
